\begin{afterword}
\summaryhead

The essay opens with Friedrich Spee, a seventeenth-century priest who described how accused witches were condemned whatever they did: a bad life proved guilt, and so did a good one; fear proved guilt, and so did calm. Having seen many cases, Spee could see that every branch led to the same verdict.

The author states a general rule and names it Conservation of Expected Evidence: before you look, the average of the probabilities you might end up with must equal the probability you hold now. If one outcome would raise your belief, the other must lower it. A likely confirmation can move you only a little; an unlikely refutation must move you a lot. So if silence is evidence for a hidden enemy, sabotage must be evidence against one, and anyone who says God hides Himself to test faith must count the miracles of the Bible against Him. A perfect reasoner cannot plan to be confirmed and can only test. The essay ends by remarking that human psychology is badly flawed.

\responsehead

The essay's mathematics is one line of a first probability course, the law of total probability. The essay gives it a new name of the author's own, modeled on the conservation laws of physics (``which I would name''). The step the name refers to, that the expected change in belief is zero, is stated in words but never written as an equation.

One paragraph, on confident theories gaining little from success and risking much from failure, states the rule in a form a reader can use. Elsewhere the essay overstates it twice. ``On average, you must expect to be exactly as confident as when you started out'' is false if confident means sure, since evidence is expected to reduce uncertainty. ``For a true Bayesian, it is impossible to seek evidence that confirms a theory'' is false in ordinary English, as the essay's own paragraph on likely confirmations shows. Both errors come from the same confusion, between ``expect'' as an average and ``expect'' in its everyday sense.

The examples show where the essay is aimed. Spee's witch trials show a tribunal that treated every outcome as guilt. Spee saw this without probability theory, and the essay does not show what the theory adds. The one example the essay builds itself is aimed at religion, and it misses. The believer's argument, as the essay states it, explains silence; it does not claim silence as evidence.

The essay ends on a joke about human psychology and gives no advice for a human judge, scientist or believer on how to act on the rule.

In short: a textbook identity given a name taken from physics, overstated twice, and aimed, in the one example the author built, at a claim the believer did not make.
\end{afterword}
